% Section 5.4, from lectures/L05/appendix/b_exercises.md. Both sets are code; their solution is in
% the repository.

\section{Exercises}
\label{c5:sec:exercises}\label{c5:app:b}

\begin{aside}
\textbf{How to work these.} Both sets extend the \code{ml} codebase from
\secref{c4:sec:exercises}: carry it forward into the lecture's exercises directory,
\file{lectures/L05/exercises} in the companion repository. Set~1 declares the
\code{Dense} class, and Set~2 implements its methods, so your network finally runs on a layer that
computes something rather than the stub it's been wired to since Chapter~\ref{c3:ch:nn}. See
\secref{c5:sec:architecture} for a walkthrough of the dense layer's architecture before you start.

\textbf{How to check your work.} The final test suite of Part~I is already in
\file{lectures/L05/exercises/test}, and Exercise~\ref{c5:ex:16} runs it. It's cumulative:
it carries the stub and network tests from Chapter~\ref{c4:ch:network} over unchanged and adds unit
tests for \code{Dense} and for the helpers in \file{ml/utils}. It tests the directory above it by
default; \code{make ML_DIR=<path>} points it at a tree of your own elsewhere.

The solution is published in the companion repository after the lecture, and
\file{lectures/L05/README.md} links to it. Try each exercise yourself before looking at it.
\end{aside}

% ------------------------------------------------------------------------------------------
\exerciseset{Declaring the Dense Layer}
\label{c5:set:1}
You'll extend your \code{ml} codebase from Chapter~\ref{c4:ch:network} by implementing a concrete
subclass \code{Dense}, which takes over from \code{ml::dense_layer::Stub} as the layer your network
is built from. The stub itself stays in the codebase; see Exercise~\ref{c5:ex:15}.

\codeexercise{The \code{ActFunc} enum class}
\label{c5:ex:1}
In \file{ml/types.hpp}, define an enum class named \code{ActFunc}. This enum class should be usable
to select the activation function for a given dense layer. Add the following enumerators:
\begin{itemize}
  \item \code{Relu}: for ReLU (\emph{Rectified Linear Unit}), which for a given input $x$ returns
    $x$ if $x > 0$, otherwise 0.
  \item \code{Tanh}: for hyperbolic tangent, which produces output in the range $(-1, 1)$.
  \item \code{None}: for no activation at all, i.e.\ the identity function, which returns $x$
    unchanged. Useful on an output layer that has to produce values outside the range a bounded
    activation function can reach.
\end{itemize}

\codeexercise{Create \file{dense.hpp}}
\label{c5:ex:2}
Create a new file named \file{dense.hpp} in \file{include/ml/dense_layer}. The file path should
therefore be \file{include/ml/dense_layer/dense.hpp}, and the file should be includable as shown
below:

\begin{cppcode}
#include "ml/dense_layer/dense.hpp"
\end{cppcode}

\codeexercise{Create \file{dense.cpp}}
\label{c5:ex:3}
Create a new file named \file{dense.cpp} in \file{source/ml/dense_layer}. The file path should
therefore be \file{source/ml/dense_layer/dense.cpp}.

\textbf{Note! Don't forget to add this file to your makefile, see below!}

\begin{makecode}
# Application target.
TARGET := app

# Source files.
# File added to the build here.
SRC_FILES := source/main.cpp\
                source/ml/dense_layer/dense.cpp\
                source/ml/neural_network/shallow.cpp\
                source/ml/utils.cpp\

# Include directory.
INC_DIR := include

# Compiler flags.
COMPILER_FLAGS := -Wall -Werror -std=c++17

# Build and run the application as default.
default: build run

# Build the application.
build:
	@g++ $(SRC_FILES) -o $(TARGET) -I$(INC_DIR) $(COMPILER_FLAGS)

# Run the application.
run:
	@./$(TARGET)

# Clean the application.
clean:
	@rm -f $(TARGET)
\end{makecode}

\codeexercise{The \code{Dense} class: declaration}
\label{c5:ex:4}
In the header file \file{include/ml/dense_layer/dense.hpp}, in the namespace
\code{ml::dense_layer}, add a class named \code{Dense} that inherits the corresponding interface,
see the file \file{include/ml/dense_layer/interface.hpp}:
\begin{itemize}
  \item Use public inheritance and mark the class \code{final} so it can't be inherited further.
  \item Override every method from the interface, including the destructor.
  \item That includes \code{initParams()}, which \code{Stub} implements as an empty method for want
    of anything to reset (see Chapter~\ref{c3:ch:nn}). \code{Dense} has bias values and weights to
    draw again, so its override does real work; the implementation follows in
    Exercise~\ref{c5:ex:10}.
\end{itemize}
\textbf{Tip:} copy the entire contents of the interface and paste it into the new file. Then adapt
the code for the new subclass \code{Dense} (no \code{virtual} or \code{= 0}, use \code{override},
etc.).

\codeexercise{Removed constructors and operators}
\label{c5:ex:5}
Delete the following:
\begin{itemize}
  \item The default constructor.
  \item The copy constructor.
  \item The move constructor.
  \item The copy operator.
  \item The move operator.
\end{itemize}

\codeexercise{Constructor}
\label{c5:ex:6}
Create a constructor that can be used to create a dense layer with arbitrary dimensions and a
selectable activation function.
\begin{itemize}
  \item \textbf{Takes:}
  \begin{itemize}
    \item \code{nodeCount}: the number of nodes in the layer (unsigned integer).
    \item \code{weightCount}: the number of weights per node in the layer (unsigned integer).
    \item \code{actFunc}: the activation function to use (of type \code{ActFunc}).
      \code{ActFunc::Relu} should be used as the default.
  \end{itemize}
  \item Should be marked \code{explicit} and \code{noexcept}.
  \item If \code{nodeCount} or \code{weightCount} equals 0, an error message should be printed and
    the program terminated by calling \code{std::terminate()}.
\end{itemize}

\codeexercise{Private member variables}
\label{c5:ex:7}
Add the following private member variables to the class:
\begin{itemize}
  \item \code{myOutput}: vector holding the nodes' output (floating-point). Should hold
    \code{nodeCount} values.
  \item \code{myPreActivationOutput}: vector holding the nodes' weighted sum before the activation
    function has been applied (floating-point). Should hold \code{nodeCount} values. Used by
    \code{backpropagate()} (implemented in Set~2) to correctly compute the activation function's
    derivative.
  \item \code{myError}: vector holding the nodes' error (floating-point). Should hold
    \code{nodeCount} values.
  \item \code{myBias}: vector holding the nodes' bias values (floating-point). Should hold
    \code{nodeCount} values.
  \item \code{myWeights}: two-dimensional vector holding the nodes' weights. Should hold
    \code{nodeCount} $\times$ \code{weightCount} values.
  \item \code{myActFunc}: the layer's activation function (of type \code{ActFunc}).
\end{itemize}

\codeexercise{Defining the methods}
\label{c5:ex:8}
Define all methods, constructors, etc.\ that aren't marked \code{delete} or \code{default} in the
file \file{source/ml/dense_layer/dense.cpp}.

\task{Constructor}
\begin{itemize}
  \item In the constructor, all member variables should be initialized:
  \begin{itemize}
    \item \code{myOutput}, \code{myPreActivationOutput}, and \code{myError} should hold
      \code{nodeCount} floating-point values equal to 0.0 at the start.
    \item \code{myBias} should hold \code{nodeCount} floating-point values randomized in
      $[0.0, 1.0]$.
    \item \code{myWeights} should hold \code{nodeCount} $\times$ \code{weightCount} floating-point
      values randomized in $[0.0, 1.0]$.
    \item The randomizing itself is written once, in \code{initParams()}, and called from here; see
      Exercise~\ref{c5:ex:10}.
    \item \code{myActFunc} should be assigned the given activation function.
  \end{itemize}
  \item If \code{nodeCount} or \code{weightCount} equals 0, the error message below should be
    printed to standard error, after which the program should terminate by calling
    \code{std::terminate()}, as shown in the code after it:
\end{itemize}

\begin{console}
Invalid dense layer parameters: nodeCount and weightCount must be greater than 0!
\end{console}

\begin{cppcode}
#include <cstdio>
#include <exception>

if ((0U == nodeCount) || (0U == weightCount))
{
    std::fprintf(stderr, "Invalid dense layer parameters: "
        "nodeCount and weightCount must be greater than 0!\n");
    std::terminate();
}
\end{cppcode}

\task{Other methods}
\begin{itemize}
  \item Implement every method from the interface except \code{feedforward()},
    \code{backpropagate()}, \code{optimize()}, and \code{initParams()}.
  \item Follow the descriptions in \file{include/ml/dense_layer/interface.hpp}:
  \begin{itemize}
    \item Methods such as \code{nodeCount()} and \code{weightCount()} should return the number of
      nodes and weights per node in the layer, respectively.
    \item Getter methods such as \code{output()} and \code{error()} should return references to the
      corresponding member variables.
    \item \code{feedforward()}, both overloads of \code{backpropagate()}, and \code{optimize()}
      return \code{bool} (see Chapter~\ref{c3:ch:nn}), so an empty body won't compile. Give each
      one a placeholder body of \code{return false;} until you implement it in Set~2.
      \code{false} rather than \code{true}, so that a method you forget to finish reports failure
      instead of quietly claiming success.
    \item \code{initParams()} returns nothing, so an empty body compiles fine. Leave it empty until
      Exercise~\ref{c5:ex:10}, where the constructor starts calling it. That's the body the stub
      keeps for good; for \code{Dense} it's a placeholder.
  \end{itemize}
\end{itemize}

% ------------------------------------------------------------------------------------------
\exerciseset{Implementing the Dense Layer}
\label{c5:set:2}
You'll complete the class \code{Dense} you just declared by implementing the methods
\code{initParams()}, \code{feedforward()}, \code{backpropagate()}, and \code{optimize()}. See
\secref{c5:sec:math} for an overview of how the math from Chapter~\ref{c3:ch:nn} maps onto the
code below.

\codeexercise{Helper functions}
\label{c5:ex:9}
\code{Dense} needs three helpers that don't belong to any one class. Add them to
\file{ml/utils.hpp} and \file{source/ml/utils.cpp}, the pair you created in
Chapter~\ref{c4:ch:network}, in the namespace \code{ml}. Declare each one in the header and define
it in the \file{.cpp}, alongside the \code{initRandGen()} already there.
\begin{itemize}
  \item \code{randomStartVal()}: function to generate and return a random number in the range
    $[0.0, 1.0]$.
  \begin{itemize}
    \item \textbf{Implementation:}
    \begin{itemize}
      \item Generate a random number in the range $[0.0, 1.0]$ by calling \code{std::rand()}, which
        generates a random number in \code{[0, RAND_MAX]}, and dividing by \code{RAND_MAX}. Both
        ranges are closed at each end: a draw of exactly \code{RAND_MAX} gives exactly \code{1.0}.
      \item One of the operands must be converted to a floating-point number to avoid integer
        division, e.g.\ \code{static_cast<double>(RAND_MAX)}.
    \end{itemize}
    \item Should be marked \code{[[nodiscard]]} and \code{noexcept}.
    \item It belongs next to \code{initRandGen()}: both wrap \code{std::rand()}, and calling this
      one without having called that one first gives the same sequence on every run.
  \end{itemize}
  \item \code{actFuncOutput()}: function to compute and return the output (floating-point) of a
    given activation function.
  \begin{itemize}
    \item \textbf{Takes:}
    \begin{itemize}
      \item \code{actFunc}: the activation function to use (of type \code{ActFunc}).
      \item \code{value}: input value to the activation function (floating-point).
    \end{itemize}
    \item \textbf{Implementation:}
    \begin{itemize}
      \item Use a switch statement to compute the output depending on the given activation
        function:
      \begin{itemize}
        \item \code{ActFunc::Relu}: return \code{value} if \code{value > 0.0}, otherwise
          \code{0.0}.
        \item \code{ActFunc::Tanh}: return \code{std::tanh(value)} (requires
          \code{#include <cmath>}).
        \item Default case (\code{ActFunc::None}): return \code{value} unchanged, which is the
          identity function.
      \end{itemize}
    \end{itemize}
    \item Should be marked \code{[[nodiscard]]} and \code{noexcept}.
  \end{itemize}
  \item \code{actFuncDelta()}: function to compute and return the derivative (floating-point) of a
    given activation function.
  \begin{itemize}
    \item \textbf{Takes:}
    \begin{itemize}
      \item \code{actFunc}: the activation function to use (of type \code{ActFunc}).
      \item \code{value}: input value to the activation function (floating-point).
    \end{itemize}
    \item \textbf{Implementation:}
    \begin{itemize}
      \item Use a switch statement to compute the derivative depending on the given activation
        function:
      \begin{itemize}
        \item \code{ActFunc::Relu}: return \code{1.0} if \code{value > 0.0}, otherwise
          \code{0.0}.
        \item \code{ActFunc::Tanh}: compute \code{const auto tanhOutput = std::tanh(value)} and
          return \code{1.0 - tanhOutput * tanhOutput}.
        \item Default case (\code{ActFunc::None}): return \code{1.0}, the derivative of the
          identity function.
      \end{itemize}
    \end{itemize}
    \item Should be marked \code{[[nodiscard]]} and \code{noexcept}.
  \end{itemize}
\end{itemize}
Neither activation function needs an error case. \code{ActFunc} has exactly three enumerators, and
all three are handled, so the default case is \code{None} rather than something that shouldn't
happen.

\codeexercise{Randomizing bias and weights}
\label{c5:ex:10}
Randomize all bias values and weights in \code{initParams()}, the interface method \code{Dense}
overrides (see Exercise~\ref{c5:ex:4}). It takes no arguments, returns nothing, and should be
marked \code{override} and \code{noexcept}:
\begin{itemize}
  \item Iterate through every node in the layer with a for loop:
    \code{for (std::size_t i{}; i < nodeCount(); ++i)}.
  \item For each node \code{i}:
  \begin{itemize}
    \item Assign a random number to its bias value by calling \code{randomStartVal()}:
      \code{myBias[i] = randomStartVal()}.
    \item Iterate through the node's weights with a nested for loop:
      \code{for (std::size_t j{}; j < weightCount(); ++j)}.
    \item For each weight \code{j}:
    \begin{itemize}
      \item Assign a random number to the weight by calling \code{randomStartVal()}:
        \code{myWeights[i][j] = randomStartVal()}.
    \end{itemize}
  \end{itemize}
\end{itemize}
In the constructor, first call \code{initRandGen()} (from \file{ml/utils.hpp}) to seed the random
number generator, then call \code{initParams()} to fill the two containers. A layer therefore
starts out randomized exactly as before; the randomization has simply moved into a method the
network can call again.

\textbf{Who calls it, and why the network needs it.} \code{Shallow::train()} calls
\code{initParams()} on both its layers before the first epoch (see Chapter~\ref{c4:ch:network}), so
every call trains a network from scratch rather than continuing the one the previous call left
behind. That matters most when a run ends badly: training on top of a network that has settled
into a bad set of parameters mostly keeps it there, while a fresh start draws new ones and can
converge on the second attempt. Without a method on the interface, the network couldn't do it at
all, since the bias and weights are private to the layer.

\codeexercise{The method \code{feedforward()}}
\label{c5:ex:11}
\task{Input validation}
\begin{itemize}
  \item Check that the dimensions of the given input match the number of weights per node in the
    layer (\code{input.size() == weightCount()}).
  \item If the dimensions don't match: print the error message
    \code{"Input dimension mismatch: expected X, actual: Y!"} and return \code{false} without
    computing anything.
\end{itemize}

\task{Computation for each node}
\begin{itemize}
  \item Iterate through every node in the layer with a for loop:
    \code{for (std::size_t i{}; i < nodeCount(); ++i)}.
  \item For each node \code{i}, compute the weighted sum:
  \begin{enumerate}
    \item Start with the node's bias value: \code{auto sum{myBias[i]}}.
    \item Add each weight multiplied by the corresponding input:
      \code{for (std::size_t j{}; j < weightCount(); ++j)} where
      \code{sum += myWeights[i][j] * input[j]}.
  \end{enumerate}
  \item Store the weighted sum before the activation function is applied:\newline
    \code{myPreActivationOutput[i] = sum}. This value is needed by \code{backpropagate()} below to
    correctly compute the activation function's derivative.
  \item Apply the activation function to the sum:
    \code{myOutput[i] = actFuncOutput(myActFunc, sum)}.
\end{itemize}

\task{Return value}
\code{true} once every node has been computed.

\codeexercise{The method \code{backpropagate()} (output layer)}
\label{c5:ex:12}
Implement \code{backpropagate()} for output layers (with reference values).

\task{Input validation}
\begin{itemize}
  \item Check that the size of the reference vector matches the number of nodes
    (\code{reference.size() == nodeCount()}).
  \item If the dimensions don't match: print the error message
    \code{"Output dimension mismatch: expected X, actual: Y!"} and return \code{false} without
    computing anything.
\end{itemize}

\task{Error computation for each node}
\begin{itemize}
  \item Iterate through every node in the layer:
    \code{for (std::size_t i{}; i < nodeCount(); ++i)}.
  \item For each node \code{i}:
  \begin{enumerate}
    \item Compute the raw error: \code{const auto err{reference[i] - myOutput[i]}}.
    \item Compute the gradient error:\newline
      \code{myError[i] = err * actFuncDelta(myActFunc, myPreActivationOutput[i])}.

      \note[Note!] Use \code{myPreActivationOutput[i]} (the weighted sum before the activation
      function was applied in \code{feedforward()}), not \code{myOutput[i]}.
      \code{actFuncDelta()} expects the activation function's \emph{input}, not its output;
      otherwise the derivative comes out wrong for \code{ActFunc::Tanh} (it happens to work for
      \code{ActFunc::Relu} by coincidence).
  \end{enumerate}
\end{itemize}

\task{Return value}
\code{true} once every node's error has been computed.

\codeexercise{The method \code{backpropagate()} (hidden layer)}
\label{c5:ex:13}
Implement \code{backpropagate()} for hidden layers (with error and weights from the next layer).

\task{Input validation}
\begin{itemize}
  \item Check that the next layer's weight count matches this layer's node count
    (\code{nextLayer.weightCount() == nodeCount()}).
  \item If the dimensions don't match: print the error message
    \code{"Layer dimension mismatch: expected X, actual: Y!"} and return \code{false} without
    computing anything.
\end{itemize}

\task{Error computation for each node}
\begin{itemize}
  \item Iterate through every node in this layer:
    \code{for (std::size_t i{}; i < nodeCount(); ++i)}.
  \item For each node \code{i}:
  \begin{enumerate}
    \item Initialize a variable to store the computed raw error: \code{double err{}}.
    \item Sum all errors from the next layer:\newline
      \code{for (std::size_t j{}; j < nextLayer.nodeCount(); ++j)}
    \begin{itemize}
      \item \code{err += nextLayer.error()[j] * nextLayer.weights()[j][i]}.
    \end{itemize}
    \item Compute the gradient error:\newline
      \code{myError[i] = err * actFuncDelta(myActFunc, myPreActivationOutput[i])}\newline (see the note in
      Exercise~\ref{c5:ex:12} for why \code{myPreActivationOutput[i]} is used instead of
      \code{myOutput[i]}).
  \end{enumerate}
\end{itemize}

\task{Return value}
\code{true} once every node's error has been computed.

\codeexercise{The method \code{optimize()}}
\label{c5:ex:14}
\task{Input validation}
\begin{itemize}
  \item Check that the learning rate lies inside the range \code{(0.0, 1.0)}, the same range the
    stub checked in Chapter~\ref{c3:ch:nn}.
  \item If the learning rate is invalid: print the error message \code{"Invalid learning rate X!"}
    and return \code{false} without updating anything.
  \item Check that the input size matches the number of weights per node
    (\code{input.size() == weightCount()}).
  \item If the dimensions don't match: print the error message
    \code{"Input dimension mismatch: expected X, actual: Y!"} and return \code{false} without
    updating anything.
\end{itemize}

\task{Parameter update for each node}
\begin{itemize}
  \item Iterate through every node in the layer:
    \code{for (std::size_t i{}; i < nodeCount(); ++i)}.
  \item For each node \code{i}:
  \begin{enumerate}
    \item Update the bias: \code{myBias[i] += myError[i] * learningRate}.
    \item Update all weights: \code{for (std::size_t j{}; j < weightCount(); ++j)}
    \begin{itemize}
      \item \code{myWeights[i][j] += myError[i] * learningRate * input[j]}
    \end{itemize}
  \end{enumerate}
\end{itemize}

\task{Return value}
\code{true} once every node's bias and weights have been updated.

\codeexercise{Compiling and testing}
\label{c5:ex:15}
Update \file{main.cpp} to use your new layer:
\begin{itemize}
  \item Replace the two \code{ml::dense_layer::Stub} instances with \code{ml::dense_layer::Dense},
    keeping the same wiring: the output layer's weight count must still match the hidden layer's
    node count.
  \item Give both layers \code{ActFunc::Tanh}, and make the hidden layer wider than the task
    strictly needs, e.g.\ 8 nodes. Keep the 2-bit XOR training data from
    Chapter~\ref{c4:ch:network}.
  \item Train for a good few thousand epochs (e.g.\ 10000), then predict once per training set and
    print the results as before. There's no learning rate to pass any more: the network sets its
    own, and stops as soon as it reaches the precision threshold, so the epoch count is an upper
    bound rather than a number to tune.
\end{itemize}
\textbf{Don't delete \file{stub.hpp}.} \code{Dense} replaces the stub in your \emph{network}, not
in your codebase. The test suite still compiles both, and still uses the stub to test
\code{Shallow}, since a layer that computes nothing is what makes the network's wiring visible. See
Exercise~\ref{c5:ex:16}.

Compile and test-run the program. Unlike Chapter~\ref{c4:ch:network}, where every prediction was
the stub's fixed output, each prediction should now land within a few hundredths of its training
target. Your exact numbers will differ from run to run, because the bias and weights start
randomized.

The early-stop line from \code{train()} should appear above the predictions, reporting a precision
and an epoch count well short of the 10000 you asked for. Both vary from run to run: the starting
parameters are random, and so is the training order. A run that stops after a few hundred epochs
and one that runs several thousand are both normal.

\textbf{When a run doesn't converge.} Roughly two XOR networks in a thousand settle into a set of
parameters they can't improve on, predicting something near 0.5 for every input however long they
train. Call \code{train()} again: it resets both layers first, so the second call starts from newly
drawn parameters rather than the ones that got stuck. That's the reset from
Chapter~\ref{c4:ch:network} earning its keep, and it's what the test suite's convergence test relies
on.

\textbf{Worth trying:} switch both layers to \code{ActFunc::Relu} and shrink the hidden layer to 2
nodes. Roughly one run in four then fails to learn XOR at all, however long you train.
\code{randomStartVal()} never returns a negative number, so every node starts on ReLU's positive
side, and a node whose weighted sum is driven negative for every input has a derivative of exactly
zero from then on: it stops learning and can't recover. That's the dying ReLU problem from
Chapter~\ref{c3:ch:nn}, and it's why the test suite trains with \code{Tanh}.

\codeexercise{Running the tests}
\label{c5:ex:16}
The final test suite is available in \file{lectures/L05/exercises/test}. It's cumulative:
it carries the stub and network tests from Chapter~\ref{c4:ch:network} over unchanged and adds unit
tests for \code{Dense} and for the helpers in \file{ml/utils}.

Once you've carried your code forward into this lecture's exercises directory, build and run it:

\begin{shell}
make -C test
\end{shell}

\noindent All 57 test cases should pass. Because \code{Dense} randomizes its own bias and weights,
none of its output can be predicted from the constructor arguments, so the tests recover the bias
by feeding the layer a vector of zeros: every weight term drops out and the weighted sum is the
bias alone. That makes every expected value exact rather than approximate.

\code{InitParamsDrawsNewValues} is the first test in the course where \code{initParams()} has
anything to do. The stub tests only check that calling it changes nothing, which is all an empty
override can be asked for; here the values it writes are read straight back out of
\code{weights()} and out of a zero-input feedforward.

The one to read is \code{BackpropagateUsesPreActivationDerivative}. It checks the note in
Exercises~\ref{c5:ex:12} and~\ref{c5:ex:13}, and only \code{ActFunc::Tanh} can catch that mistake:
for \code{Relu} the two values agree, since \code{max(0.0, s)} is positive exactly when \code{s}
is.

The test suite's README, \file{lectures/L05/exercises/test/README.md}, has more
information, including how the convergence threshold was chosen.
